\subsection*{Regional cohort literacy after 1939}

Regional tabulations of the 1959 census allow a more demanding follow-up
than the national comparison alone. I compare the cohorts aged 30--34,
35--39 and 40--44 in January 1939 with those aged 50--54, 55--59 and 60--64
in January 1959, holding their age composition fixed at the 1939 census.
These age bands can be matched without interpolation. They overlap with,
but do not exactly reproduce, the paper's group aged 20--34 in December
1926. The five-republic diagnostic covers the Tatar, Bashkir, Chuvash,
Mari and Udmurt republics. Karelia is excluded because of territorial
comparability concerns. The correspondence follows the republics'
territorial cores; a complete concordance of small territorial transfers
has not been established.

Across these five republics, female literacy in the selected cohorts
increases by \SN{followup.total.F} percentage points for the total population and by \SN{followup.rural.F}
points for the rural population, but changes by \SN{followup.urban.F} points in towns.
The corresponding changes for men are \SN{followup.total.M}, \SN{followup.rural.M} and \SN{followup.urban.M} points.
An exploratory regression of the cohort-level change on a female
indicator, republic and age-group fixed effects, weighted by population
in 1939, gives a female-minus-male change of \SN{urban.b} points in towns.
The CR2/Satterthwaite 95 per cent interval is [\SN{urban.ci_low_cr2}, \SN{urban.ci_high_cr2}]; the
restricted wild cluster bootstrap, enumerating all 32 Rademacher sign
vectors for the five republics, gives $p=\SN{urban.wcr_p}$. The two procedures give different answers at the conventional five per cent threshold, so both are reported.

These are changes in the recorded literacy of populations occupying
successive age groups, not transitions observed for linked individuals.
For example, the Tatar republic's female urban cohort population is
120.9 per cent of its 1939 size by 1959. Migration, settlement
reclassification, territorial change or enumeration error must therefore
enter the comparison. Wartime mortality, selective survival, subsequent
learning and possible loss of skills add further components.
Consequently, similarity between the two national literacy rates cannot
by itself bound misreporting in 1939 or establish uniform regional
persistence. The post-1939 gender-gap estimates are descriptive and do
not identify the causal effect of the literacy campaign.

% Source URLs and file hashes are recorded in the source register and manifests.
